\documentclass[12pt]{article}
\usepackage[utf8]{inputenc}
\usepackage[T1]{fontenc}
\usepackage{amsmath, amssymb, geometry, booktabs}
\usepackage{enumitem}
\geometry{margin=1in}
\setlength{\parindent}{0pt}
\renewcommand{\arraystretch}{1.2}

\begin{document}


\section*{Problem 2 (13 marks)}

A country produces \textbf{agricultural products (A)} and \textbf{electronics (E)} with the production possibility frontier (PPF):
\[
A^2 + 2E^2 = 900.
\]

The world relative price is \( P_E / P_A = 3 \).  
The country consumes at the point where preferences satisfy \( E = 0.5A \).

\begin{enumerate}[label=(\arabic*)]
\item Determine the production point $(A_P, E_P)$ that maximizes the value of output at world prices.  
\item Find the consumption bundle $(A_C, E_C)$, identify which good is exported, and compute the value of exports and imports.  
\item Verify that trade is balanced at world prices.  
\item Explain in words how an increase in \( P_E / P_A \) (improvement in the terms of trade) affects the country’s production, consumption, and welfare.  
\item Suppose the economy experiences \textbf{export-biased growth} that expands production possibilities more in electronics than in agriculture.  
\begin{itemize}
    \item Describe how this changes the country’s trade pattern and the direction of the terms-of-trade effect.  
    \item Discuss whether welfare necessarily increases after such growth.
\end{itemize}
\end{enumerate}



\end{document}
